\documentclass[10pt,a4paper]{amsart}
\usepackage{amsmath,amssymb,amsthm,graphicx}
% Retained minimal preamble for the Stacks Project chapter more-algebra.tex (More on Algebra).
% The chapter file pinned-original-source.tex carries only \input{preamble} on its first line, so no
% preamble text lives inside the chapter itself.  The Stacks Project uses one single shared preamble for
% every chapter; the snapshot stages it identically as sources/stacks-chow/preamble.tex and
% sources/stacks-etale-cohomology/preamble.tex.  The block below is copied mechanically, line for line,
% from sources/stacks-chow/preamble.tex lines 14-18 (the Xy-pic setup) and 154-217 (the native theorem
% declarations together with the mathematical macros).  Nothing was added, removed or redefined; the
% wrapper class and the amsmath/amssymb/amsthm/graphicx base packages are supplied by the assembler.
% Chapter: more-algebra.tex -- More on Algebra.
% For commutative diagrams we use Xy-pic
\usepackage[all]{xy}

% We use 2cell for 2-commutative diagrams.
\xyoption{2cell}

% Theorem environments.
%
\theoremstyle{plain}
\newtheorem{theorem}[subsection]{Theorem}
\newtheorem{proposition}[subsection]{Proposition}
\newtheorem{lemma}[subsection]{Lemma}

\theoremstyle{definition}
\newtheorem{definition}[subsection]{Definition}
\newtheorem{example}[subsection]{Example}
\newtheorem{exercise}[subsection]{Exercise}
\newtheorem{situation}[subsection]{Situation}

\theoremstyle{remark}
\newtheorem{remark}[subsection]{Remark}
\newtheorem{remarks}[subsection]{Remarks}

\numberwithin{equation}{subsection}

% Macros
%
\def\lim{\mathop{\mathrm{lim}}\nolimits}
\def\colim{\mathop{\mathrm{colim}}\nolimits}
\def\Spec{\mathop{\mathrm{Spec}}}
\def\Hom{\mathop{\mathrm{Hom}}\nolimits}
\def\Ext{\mathop{\mathrm{Ext}}\nolimits}
\def\SheafHom{\mathop{\mathcal{H}\!\mathit{om}}\nolimits}
\def\SheafExt{\mathop{\mathcal{E}\!\mathit{xt}}\nolimits}
\def\Sch{\mathit{Sch}}
\def\Mor{\mathop{\mathrm{Mor}}\nolimits}
\def\Ob{\mathop{\mathrm{Ob}}\nolimits}
\def\Sh{\mathop{\mathit{Sh}}\nolimits}
\def\NL{\mathop{N\!L}\nolimits}
\def\CH{\mathop{\mathrm{CH}}\nolimits}
\def\proetale{{pro\text{-}\acute{e}tale}}
\def\etale{{\acute{e}tale}}
\def\QCoh{\mathit{QCoh}}
\def\Ker{\mathop{\mathrm{Ker}}}
\def\Im{\mathop{\mathrm{Im}}}
\def\Coker{\mathop{\mathrm{Coker}}}
\def\Coim{\mathop{\mathrm{Coim}}}

% Boxtimes
%
\DeclareMathSymbol{\boxtimes}{\mathbin}{AMSa}{"02}

%
% Macros for moduli stacks/spaces
%
\def\QCohstack{\mathcal{QC}\!\mathit{oh}}
\def\Cohstack{\mathcal{C}\!\mathit{oh}}
\def\Spacesstack{\mathcal{S}\!\mathit{paces}}
\def\Quotfunctor{\mathrm{Quot}}
\def\Hilbfunctor{\mathrm{Hilb}}
\def\Curvesstack{\mathcal{C}\!\mathit{urves}}
\def\Polarizedstack{\mathcal{P}\!\mathit{olarized}}
\def\Complexesstack{\mathcal{C}\!\mathit{omplexes}}
% \Pic is the operator that assigns to X its picard group, usage \Pic(X)
% \Picardstack_{X/B} denotes the Picard stack of X over B
% \Picardfunctor_{X/B} denotes the Picard functor of X over B
\def\Pic{\mathop{\mathrm{Pic}}\nolimits}
\def\Picardstack{\mathcal{P}\!\mathit{ic}}
\def\Picardfunctor{\mathrm{Pic}}
\def\Deformationcategory{\mathcal{D}\!\mathit{ef}}
\title{Power torsion depends only on the closed subset: a module is I-power torsion exactly when its support lies in V(I)}
\author{The Stacks Project Authors}
\date{stacks-project:more-algebra.tex:local-snapshot-20261012; GNU FDL 1.2 or later}
\begin{document}
\maketitle
\setcounter{subsection}{1}
\noindent\textit{Editorial note.} All delivered mathematical text is the author's own, copied from the pinned Stacks Project chapter file more-algebra.tex; no statement or proof step was rewritten, reordered or removed. The counted statement is the labeled lemma environment at source lines 24085-24092 (canonical tag 0953) and the counted proof is the proof environment at source lines 24094-24104, which immediately follows it and is its own proof.\par\medskip
\noindent\textit{Editorial note.} The proof uses the macros Spec, Ker and the module categories of the shared preamble, all of which the retained preamble supplies verbatim.\par\medskip
\noindent\textit{Editorial note.} The sentence at source line 24106, remarking that the next two lemmas should probably go somewhere else, lies just after the counted proof and is outside the counted spans.\par\medskip
\noindent\textit{Editorial note.} The chapter file pinned-original-source.tex carries only an input of the shared preamble on its first line, so no preamble text lives inside the chapter. The retained file sources/stacks-more-algebra/retained-preamble.tex is a mechanical, line-for-line copy of the project's single shared preamble as staged in the snapshot (Xy-pic setup plus the native theorem declarations and mathematical macros). Nothing was added, removed or redefined, and the excerpt serves only as preamble, never as a span coordinate.\par\medskip
\noindent\textit{Editorial note.} The source URL is the source-level Stacks Project URL, as required when more than one card is produced from one source; the per-claim tag is carried by claim\_key and by this note. The counted statement and proof refer to no labelled item of the chapter and cite no bibliography entry, so this card contributes no external reference stub and needs no bibliography.\par\medskip
\noindent\textit{Allowed standard prerequisites (editorial).}\begin{itemize}
\item An element of a module that maps to zero in the localization at f is annihilated by a power of f, and vanishing of the localizations at every prime not containing I is equivalent to vanishing of the localizations at the elements f of I when I is finitely generated. used at proof lines 24095-24100, where the passage from localizations at primes to localizations at the generators f\_1, ..., f\_r of I is made. Converts the support condition into explicit annihilation by the generators of I.
\item The support of an R-module is the set of primes at which its localization is nonzero, and V(I) is the set of primes containing I. used throughout statement and proof, notably at proof line 24097 and at the description of the kernel at proof lines 24101-24102. Provides the vocabulary in which both the hypothesis and the conclusion of the lemma are stated.
\item The submodule of elements annihilated by a power of I is the kernel of the map to the product of the localizations at the primes outside V(I). used at proof lines 24101-24103. Turns the elementwise equivalence into an identification of the I-power torsion submodule, from which both assertions follow.
\end{itemize}\medskip
\begin{lemma}
\label{lemma-local-cohomology-closed}
Let $R$ be a ring and let $I \subset R$ be a finitely generated ideal.
The subcategory $I^\infty\text{-torsion} \subset \text{Mod}_R$
depends only on the closed subset $Z = V(I) \subset \Spec(R)$.
In fact, an $R$-module $M$ is $I$-power torsion if and only if its
support is contained in $Z$.
\end{lemma}

\begin{proof}
Let $M$ be an $R$-module. Let $x \in M$. If $x \in M[I^\infty]$, then $x$
maps to zero in $M_f$ for all $f \in I$. Hence $x$ maps to zero in
$M_\mathfrak p$ for all $\mathfrak p \not \supset I$. Conversely, if $x$
maps to zero in $M_\mathfrak p$ for all $\mathfrak p \not \supset I$,
then $x$ maps to zero in $M_f$ for all $f \in I$. Hence if
$I = (f_1, \ldots, f_r)$, then $f_i^{n_i}x = 0$ for some $n_i \geq 1$.
It follows that $x \in M[I^{\sum n_i}]$. Thus $M[I^\infty]$ is
the kernel of $M \to \prod_{\mathfrak p \not \in Z} M_\mathfrak p$.
The second statement of the lemma follows and it implies the first.
\end{proof}

\end{document}
